% ============================================================
% CHAPITRE 6 : Régulateurs, Correcteurs, Boucle fermée
% ============================================================
\chapterbanner{Chapitre 6 --- R\'egulateurs \& Boucle ferm\'ee}{ch6teal}

\begin{tcolorbox}[colback=ch6teal!5,colframe=ch6teal!80!black,title={Notations (rappel cours)}]
\begin{itemize}
  \item $G_c(s)$ : \textbf{r\'egulateur} (correcteur)
  \item $G_a(s)$ : \textbf{syst\`eme \`a r\'egler} (ampli + actionneur + processus + capteur) ; $K_a{=}G_a(0)$ = gain du proc\'ed\'e
  \item $G_o(s)=G_c\,G_a$ : \textbf{boucle ouverte} (not\'ee aussi $L$) ; $K_o{=}\lim_{s\to0}G_o$ = gain permanent de boucle
  \item $G_{yw}=\frac{Y}{W}=\frac{G_o}{1+G_o}$ : BF \textbf{correspondance} (poursuite)
  \item $G_{ew}=\frac{E}{W}=\frac{1}{1+G_o}$ : transfert de l'\textbf{erreur} $=S$
  \item $G_{yv}=\frac{Y}{V}=\frac{G_p}{1+G_o}$ : BF \textbf{maintien} (rejet de $v$ ; $G_p$ = partie en aval de $v$)
  \item $G_m$ : \textbf{capteur} (retour) ; r\'etroaction unitaire $\Leftrightarrow G_m{=}1$
  \item $S=G_{ew}$ (sensibilit\'e), $T=G_{yw}$ (compl\'em.) : $S+T=1$
  \item \textbf{Param. r\'egulateur} : $K_p$ (gain prop.), $T_i$ (temps int\'egral), $T_d$ (temps d\'eriv\'e), $K_i{=}K_p/T_i$ ; $\alpha$ : facteur du PD r\'ealisable ($\approx0.1$--$0.2$)
  \item \textbf{Constantes d'erreur} : $K_o$ (position), $K_v{=}\lim_{s\to0}sG_o$ (vitesse), $K_a^{\text{acc}}{=}\lim_{s\to0}s^2G_o$ (acc\'el.)
  \item $\varphi_m, A_m$ : marges de phase / gain ; $\omega_{co}$ : coupure 0~dB en BO ; $\omega_\pi$ : phase $-180°$
  \item $K_{cr}, T_{cr}$ : gain et p\'eriode critiques (limite de stabilit\'e)
  \item \textbf{Perf.} (cf. ch.3) : $D_\%$ d\'ep., $t_r$ mont\'ee, $T_{\text{r\'eg}}$ r\'eglage, $\zeta$ amort., $\omega_n$ puls. nat.
\end{itemize}
\end{tcolorbox}

\begin{tcolorbox}[colback=ch6teal!5,colframe=ch6teal!80!black,title={Transferts en boucle ferm\'ee}]
Avec $G_o(s) = G_c(s) \cdot G_a(s)$ (transfert de boucle ouverte ; $G_a$ = syst\`eme \`a r\'egler), r\'etroaction unitaire :
\tightmath{
\boxed{G_{yw} = \frac{G_o}{1{+}G_o}} \quad \boxed{G_{ew} = \frac{1}{1{+}G_o}} \quad \boxed{G_{yv} = \frac{G_p}{1{+}G_o}}
}
\piege{Le d\'enom. $1{+}G_o(s)$ est commun \`a \textbf{tous} les transferts BF. \'Eq. caract. : $1{+}G_o(s){=}0$. $S{=}G_{ew}$, $T{=}G_{yw}$, $S{+}T{=}1$ (voir bo\^ite sensibilit\'e).}
\end{tcolorbox}

\begin{tcolorbox}[colback=ch6teal!5,colframe=ch6teal!80!black,title={Erreur permanente : constantes, tableau \& perturbation}]
$K_o = \lim_{s\to 0} G_o(s)$ (position), $K_v = \lim_{s\to 0} sG_o(s)$ (vitesse), $K_a^{\text{acc}} = \lim_{s\to 0} s^2 G_o(s)$ (acc\'el.)
\begin{center}
\renewcommand{\arraystretch}{1.1}
\scriptsize
\begin{tabular}{@{}lccc@{}}
\toprule
Entr\'ee $w(t)$ & Type 0 & Type 1 & Type 2 \\
\midrule
\'Echelon $A$ & $\frac{A}{1{+}K_o}$ & $0$ & $0$ \\[3pt]
Rampe $At$ & $\infty$ & $\frac{A}{K_v}$ & $0$ \\[3pt]
Parabole $\frac{A}{2}t^2$ & $\infty$ & $\infty$ & $\frac{A}{K_a^{\text{acc}}}$ \\
\bottomrule
\end{tabular}
\end{center}
\vars{$A$ = \textbf{amplitude de la consigne $w$} (hauteur du saut, pente de la rampe, ou accel.). En Laplace : \'echelon $W{=}A/s$, rampe $W{=}A/s^2$, parabole $W{=}A/s^3$. Perturbation : amplitude $B$ ($V{=}B/s$).}
\textbf{Poursuite} ($v{=}0$) : $e_\infty = \frac{A}{1{+}K_o}$ \quad
\textbf{Maintien} ($w{=}0$) : $e_\infty = -\lim_{s\to 0} \frac{G_p(s)}{1{+}G_o(s)}$
\par\smallskip
\textbf{$w$ et $v$ simultan\'es} (superposition, $W{=}A/s$, $V{=}B/s$) :
\tightmath{
y_\infty = G_{yw}(0)\,A + G_{yv}(0)\,B \qquad e_\infty = A - y_\infty
}
\textbf{Type 0} : $G_{yw}(0){=}\frac{K_o}{1{+}K_o}$, $G_{yv}(0){=}\frac{G_p(0)}{1{+}K_o}$ \quad
\textbf{Type 1} : $G_{yw}(0){=}1$, $G_{yv}(0){=}\lim_{s\to 0}\frac{G_p(s)}{1{+}G_o(s)}$
\vars{Ex. type 1 : $G_o = 2K_p/[s(1{+}s)]$, $G_{yv}(0) = 1/(2K_p)$. Pour $V{=}0.5$ : $y_{v,\infty} = 0.5/(2K_p)$.}
\piege{Type $\alpha$ \'elev\'e $\Rightarrow$ meilleure pr\'ecision, mais plus difficile \`a stabiliser.}
\end{tcolorbox}

\begin{tcolorbox}[colback=ch6teal!5,colframe=ch6teal!80!black,title={R\'egulateur P (Proportionnel)}]
\tightmath{
\boxed{G_c(s) = K_p}
}
\begin{itemize}
  \item R\'eduit $t_r$ (augmente $\omega_{co}$), augmente $D_\%$ ($\varphi_m \downarrow$)
  \item R\'eduit mais \textbf{ne supprime pas} $e_\infty$ (type 0)
  \item \textbf{Dilemme} : $K_p \uparrow \Rightarrow e_\infty \downarrow$ mais stabilit\'e $\downarrow$
\end{itemize}
\end{tcolorbox}

\begin{tcolorbox}[colback=ch6teal!5,colframe=ch6teal!80!black,title={R\'egulateur PI}]
\tightmath{
\boxed{G_c(s) = K_p\!\left(1 + \frac{1}{T_i s}\right) = K_p\,\frac{1 + T_i s}{T_i s}}
}
\begin{itemize}
  \item Ajoute int\'egrateur $\Rightarrow$ type $+1$ $\Rightarrow$ \textbf{supprime $e_\infty$} (\'echelon)
  \item Z\'ero en $s = -1/T_i$ : compenser p\^ole dominant
  \item Diminue $\varphi_m$ (attention stabilit\'e)
\end{itemize}
\astuce{\textbf{R\'eglage PI} : 1. $T_i = \tau_{\text{proc}}$ (annule p\^ole dominant). 2. R\'egler $K_p$ pour $\varphi_m$ souhait\'ee. 3. V\'erifier $\omega_{co}$ et $A_m$.}
\end{tcolorbox}

\begin{tcolorbox}[colback=ch6teal!5,colframe=ch6teal!80!black,title={R\'egulateur PD \& PID}]
\textbf{PD} : $\boxed{G_c(s) = K_p(1 + T_d s)}$ \quad \textbf{PD r\'ealisable} (\og r\'egulateur \`a avance de phase\fg{}, facteur $\alpha\approx0.1$--$0.2$) :
\tightmath{
\boxed{G_c = K_p\!\left(1 + \frac{sT_d}{1{+}\alpha s T_d}\right) = K_p\,\frac{1{+}(1{+}\alpha)T_d s}{1{+}\alpha T_d s}}
}
\begin{itemize}
  \item Z\'ero \`a $\omega_z = \frac{1}{(1{+}\alpha)T_d}$, p\^ole \`a $\omega_p = \frac{1}{\alpha T_d}$, gain BF $= K_p$
  \item Am\'eliore amortissement, r\'eduit $D_\%$, augmente $\varphi_m$. Amplifie bruit HF.
\end{itemize}
\vars{\textbf{Identification depuis Bode} : $K_p = 10^{G_{\text{BF}}/20}$, puis z\'ero $\omega_z$ et p\^ole $\omega_p$ $\Rightarrow$ $T_d = 1/\omega_z - 1/\omega_p$, $\alpha = \omega_z/(\omega_p - \omega_z)$.}
\textbf{PID} (forme du cours) :
\tightmath{
\boxed{G_c(s) = K_p\!\left(1 + \frac{1}{T_i s} + T_d s\right) = K_p\,\frac{1 + s T_i + s^2 T_i T_d}{s T_i}}
}
Num\'erateur factoris\'e : $1{+}sT_i{+}s^2 T_i T_d = (1{+}s\tau_1)(1{+}s\tau_2)$ avec $\tau_1{+}\tau_2 = T_i$, $\tau_1\tau_2 = T_i T_d$.
\end{tcolorbox}

\begin{tcolorbox}[colback=ch6teal!5,colframe=ch6teal!80!black,title={Ziegler-Nichols, gain critique $K_{cr}$ \& p\'eriode $T_{cr}$}]
\textbf{M\'ethode BO} (r\'ep. indic.) : tangente au point d'inflexion $\Rightarrow$ retard $T_u$ et temps de mont\'ee de la tangente $T_g$ (pente $R{=}\Delta y/T_g$).
\begin{center}
\renewcommand{\arraystretch}{1}
\scriptsize
\begin{tabular}{@{}lccc@{}}
\toprule
R\'eg. & $K_p$ & $T_i$ & $T_d$ \\
\midrule
P & $T_g/T_u$ & -- & -- \\
PI & $0.9\,T_g/T_u$ & $3.3\,T_u$ & -- \\
PID & $1.2\,T_g/T_u$ & $2\,T_u$ & $0.5\,T_u$ \\
\bottomrule
\end{tabular}
\end{center}
\textbf{M\'ethode BF} (oscillation) : $K_{cr}$ = gain du P amenant la BF \`a la \textbf{limite de stabilit\'e} ; $T_{cr}$ = p\'eriode de l'oscillation. \textbf{Lecture sur le Bode de $G_o$} :
\tightmath{
\boxed{K_{cr} = \frac{1}{|G_o(\jw_\pi)|}} \qquad \boxed{T_{cr} = \frac{2\pi}{\omega_\pi}}
}
\begin{center}
\renewcommand{\arraystretch}{1}
\scriptsize
\begin{tabular}{@{}lccc@{}}
\toprule
R\'eg. & $K_p$ & $T_i$ & $T_d$ \\
\midrule
P & $0.5\,K_{cr}$ & -- & -- \\
PI & $0.45\,K_{cr}$ & $T_{cr}/1.2$ & -- \\
PID & $0.6\,K_{cr}$ & $T_{cr}/2$ & $T_{cr}/8$ \\
\bottomrule
\end{tabular}
\end{center}
\piege{Ziegler-Nichols donne $D_\% \approx 25\%$. Souvent il faut ajuster ensuite. \`A $K{=}K_{cr}$ le lieu passe exactement par $-1$.}
\end{tcolorbox}

\begin{tcolorbox}[colback=ch6teal!5,colframe=ch6teal!80!black,title={Questions qualitatives : effet d'un r\'eglage sur la BF}]
\textbf{Recette} : toute question \og effet de \dots\ ? \fg{} se traite sur \textbf{3 axes}, chacun \textbf{justifi\'e} par sa cha\^ine cause$\to$effet.
\begin{enumerate}\setlength\itemsep{0pt}
  \item \textbf{Stabilit\'e} ($\varphi_m, A_m$) : $K_p$ ne d\'eplace que le \textbf{gain} BO $\Rightarrow$ d\'eplace $\omega_{co}$ ; $\varphi_m$ se lit \`a $\omega_{co}$, $A_m$ \`a $\omega_\pi$.
  \item \textbf{Pr\'ecision} ($e_\infty$) : regarder le \textbf{TYPE} (nb d'int\'egrateurs), \emph{pas} la valeur de $K_p$. Int\'egrateur pur $\Rightarrow e_\infty^{\text{\'ech}}{=}0$ ; type 0 $\Rightarrow e_\infty{=}\frac{A}{1+K_o}$ ($K_o{\propto}K_p$).
  \item \textbf{Rapidit\'e} ($t_r, T_{\text{r\'eg}}$) : bande passante BF $\omega_{-3\text{dB}}^{BF}\approx\omega_{co}$ ; $T_{\text{r\'eg}}\approx\pi/\omega_{co}$.
\end{enumerate}
\textbf{Ex. r\'eduire $K_p$} (gain BO $\downarrow$) :
\begin{itemize}\setlength\itemsep{0pt}
  \item Stab. \textbf{$\uparrow$} : $\omega_{co}\downarrow$ $\Rightarrow$ phase plus haute $\Rightarrow \varphi_m\uparrow$ ; lieu s'\'eloigne de $-1 \Rightarrow A_m\uparrow$.
  \item Pr\'ec. \textbf{$=$} si int\'egrateur pur ($e_\infty{=}0$), sinon \textbf{$\downarrow$} ($e_\infty\uparrow$).
  \item Rapid. \textbf{$\downarrow$} : $\omega_{co}\downarrow \Rightarrow$ BF plus lente ($T_{\text{r\'eg}}\uparrow$).
\end{itemize}
\astuce{\textbf{Justifier par la cha\^ine}, jamais juste \og \c{c}a augmente \fg{} : \og $K_p\downarrow \Rightarrow$ gain BO$\downarrow \Rightarrow \omega_{co}\downarrow \Rightarrow \varphi_m\uparrow$ \fg{}. Pour $K_p\uparrow$ : inverser les fl\`eches.}
\end{tcolorbox}

\begin{tcolorbox}[colback=ch6teal!5,colframe=ch6teal!80!black,title={R\'egulateurs : comparatif, effets \& choix}]
\renewcommand{\arraystretch}{1.18}
\scriptsize
\begin{tabularx}{\linewidth}{@{}l l c >{\raggedright\arraybackslash}X@{}}
\toprule
R\'eg. & $G_c(s)$ & $e_\infty$ \'ech. & Effet principal \\
\midrule
P   & $K_p$ & $\frac{A}{1+K_o}{\neq}0$ & $t_r{\downarrow}$, $\omega_{co}{\uparrow}$, $D_\%{\uparrow}$, $\varphi_m{\downarrow}$ \\
I   & $\frac{K_i}{s}$ & $0$ & type $+1$ ; $\varphi_m{\downarrow}$ ($-90°$), lent, windup \\
PI  & $K_p\frac{1{+}T_i s}{T_i s}$ & $0$ & pr\'ecision + compense p\^ole dom. ; $\varphi_m{\downarrow}$ \\
PD  & $K_p(1{+}T_d s)$ & inchang\'e & $D_\%{\downarrow}$, $\varphi_m{\uparrow}$, $t_r{\downarrow}$ ; amplifie bruit HF \\
PID & $K_p(1{+}\tfrac{1}{T_i s}{+}T_d s)$ & $0$ & pr\'ecision \textbf{et} amortissement \\
\bottomrule
\end{tabularx}
\par\smallskip
\textbf{Effet d'une action sur la BF} :
\begin{center}
\renewcommand{\arraystretch}{1.1}
\scriptsize
\begin{tabular}{@{}lccccc@{}}
\toprule
 & $t_r$ & $D_\%$ & $e_\infty$ & $\varphi_m$ & $\omega_{co}$ \\
\midrule
$K_p \uparrow$ & $\downarrow$ & $\uparrow$ & $\downarrow$ & $\downarrow$ & $\uparrow$ \\
$+$ I ($T_i$) & -- & $\uparrow$ & $\boldsymbol{= 0}$ & $\downarrow$ & -- \\
$+$ D ($T_d$) & -- & $\downarrow$ & -- & $\uparrow$ & -- \\
\bottomrule
\end{tabular}
\end{center}
\textbf{Quel choisir ?} R\'eduire $e_\infty$ en gardant $\varphi_m$ : $K_p\uparrow$. Supprimer $e_\infty$ \'echelon : PI. Supprimer $e_\infty$ rampe : PI sur proc. type\,1. R\'eduire $D_\%$/$t_r$, augm. $\varphi_m$ : PD. Les deux : PID.
\vars{P : rapide mais erreur r\'esiduelle. I : annule $e_\infty$ mais d\'estabilise. D : anticipe/amortit, sensible au bruit. PI = pr\'ecision, PD = stabilit\'e/rapidit\'e, PID = les deux.}
\end{tcolorbox}

\begin{tcolorbox}[colback=ch6teal!5,colframe=ch6teal!80!black,title={Transferts BF avec capteur $G_m \neq 1$}]
\tightmath{
G_{yw} = \frac{G_c\,G_a}{1 + G_c\,G_a\,G_m} \quad G_o = G_c\,G_a\,G_m
}
\vars{$G_a$ = syst\`eme \`a r\'egler (cha\^ine directe), $G_m$ = capteur (retour). R\'etroaction unitaire $\Leftrightarrow G_m{=}1$.}
\end{tcolorbox}

\begin{tcolorbox}[colback=ch6teal!5,colframe=ch6teal!80!black,title={R\'esum\'e des fonctions de transfert types}]
\begin{center}
\renewcommand{\arraystretch}{1.1}
\scriptsize
\begin{tabular}{@{}lll@{}}
\toprule
Nom & $G(s)$ & P\^oles \\
\midrule
1\textsuperscript{er} ordre & $\frac{K}{1{+}\tau s}$ & $-1/\tau$ \\[2pt]
2\textsuperscript{nd} ordre & $\frac{K\omega_n^2}{s^2{+}2\zeta\omega_n s{+}\omega_n^2}$ & $-\zeta\omega_n \pm j\omega_0$ \\[2pt]
Int\'egrateur & $K/s$ & $0$ \\[2pt]
Int. + 1\textsuperscript{er} o. & $K/[s(1{+}\tau s)]$ & $0,\;-1/\tau$ \\[2pt]
Retard & $Ke^{-sT_r}$ & aucun \\[2pt]
PI & $K_p\frac{1{+}T_i s}{T_i s}$ & $0$ ; z: $-1/T_i$ \\[2pt]
PD r\'eal. & $K_p\frac{1{+}(1{+}\alpha)T_d s}{1{+}\alpha T_d s}$ & $-\frac{1}{\alpha T_d}$ ; z: $-\frac{1}{(1{+}\alpha)T_d}$ \\
\bottomrule
\end{tabular}
\end{center}
\vars{P\^ole dominant asservi : $t_{s5} \approx 3/|\text{Re}(s_{\text{dom}})|$, $t_{s2} \approx 4/|\text{Re}(s_{\text{dom}})|$ (dominant si les autres p\^oles ont $|\text{Re}|\ge 5\times$ plus grande).}
\end{tcolorbox}

\begin{tcolorbox}[colback=ch6teal!5,colframe=ch6teal!80!black,title={Sensibilit\'e \& robustesse}]
\tightmath{
\boxed{S(s) = \frac{1}{1{+}G_o},\quad T(s) = \frac{G_o}{1{+}G_o},\quad S + T = 1}
}
\begin{itemize}
  \item $S$ : passe-haut $\Rightarrow$ att\'enue les perturbations BF ; $T$ : passe-bas
  \item Robustesse : $\|S\|_\infty \leq 2$ ($6$~dB) recommand\'e $\Leftrightarrow \varphi_m \geq 30°$
  \item $\|S\|_\infty = 1/d_{\min}$, \; $d_{\min} = \min_\omega|1{+}G_o(\jw)|$ (distance critique, \S5.4.3)
\end{itemize}
\piege{Minimiser $e_\infty$ et maximiser $\varphi_m$ sont conflictuels : dilemme pr\'ecision-stabilit\'e.}
\end{tcolorbox}

\begin{tcolorbox}[colback=ch6teal!5,colframe=ch6teal!80!black,title={R\'eponse harmonique BF : $G(j\omega)$ \& calcul num\'erique}]
\textbf{Identit\'es} (substituer $s = j\omega$) :
\tightmath{
|j\omega| = \omega,\quad |1{+}j\omega\tau| = \sqrt{1{+}(\omega\tau)^2}
}
\tightmath{
\angle j\omega = 90°,\quad \angle(1{+}j\omega\tau) = \arctan(\omega\tau)
}
\vars{Produit : $|AB|=|A||B|$, $\angle AB = \angle A + \angle B$ ; quotient : $\angle\frac{A}{B}=\angle A-\angle B$. $K{>}0$ : $\angle K{=}0°$ ; $K{<}0$ : $-180°$.}
\textbf{Sortie BF} pour $w(t)=A\sin(\omega t)$ : $y(t) = A\,|G_{BF}(\jw)|\,\sin(\omega t + \angle G_{BF}(\jw))$.
\par\smallskip
\textbf{M\'ethode depuis Bode BO} :
\begin{enumerate}\setlength\itemsep{0pt}
  \item Lire $|G_o|$ et $\angle G_o$ \`a $\omega$ (pour $K_{p,0}$ du Bode)
  \item Corriger : $|G_o|_{\text{r\'eel}} = |G_o| \cdot K_p/K_{p,0}$, $\angle$ inchang\'e
  \item $1{+}G_o = (1+M\cos\theta)+jM\sin\theta$ avec $G_o = M\angle\theta$
  \item $|G_{BF}| = M/|1{+}G_o|$, \quad $\angle G_{BF} = \theta - \angle(1{+}G_o)$
\end{enumerate}
\piege{Si Re$(1{+}G_o) < 0$ : tenir compte du quadrant (atan2).}
\end{tcolorbox}

\begin{tcolorbox}[colback=ch6teal!5,colframe=ch6teal!80!black,title={Lire le Bode \`a un autre $K_p$ (d\'ecalage de gain)}]
Le Bode est trac\'e pour un $K_p^{\text{base}}$ donn\'e (souvent $1$ ou $100$). Changer $K_p$ \textbf{d\'ecale verticalement la courbe de gain}, la \textbf{phase est inchang\'ee} :
\tightmath{
\boxed{\;|G_o|^{\text{dB}}_{K_p} = |G_o|^{\text{dB}}_{K_p^{\text{base}}} + 20\log_{10}\!\frac{K_p}{K_p^{\text{base}}}\;},\qquad \angle G_o \text{ inchang\'e}
}
\textbf{Ex.} lire $K_p{=}2$ : base $100 \Rightarrow 20\log\frac{2}{100}{=}-34$ dB (descendre) ; base $1 \Rightarrow +6$ dB (monter).
\vars{Cons\'equence : $\omega_\pi$ \textbf{inchang\'ee} (phase intacte) mais $\omega_{co}$ se d\'eplace $\Rightarrow A_m$ et $\varphi_m$ changent. Rep\`eres : $\times2{\to}{+}6$ dB, $\times10{\to}{+}20$ dB, $\div10{\to}{-}20$ dB.}
\end{tcolorbox}

\begin{tcolorbox}[colback=ch6teal!5,colframe=ch6teal!80!black,title={Correspondance BO $\to$ BF : r\`egles de pouce}]
\begin{itemize}
  \item $\omega_{-3\text{dB}}^{BF} \approx \omega_{co}$ pour $30° < \varphi_m < 70°$
  \item Bode de $G_{BF}$ : pic de r\'esonance \`a $\omega_r \approx \omega_{co}$ si $\varphi_m$ petit
  \item $\zeta \approx \varphi_m\!{}^\circ/100$ puis $D_\% = 100\,e^{-\pi\zeta/\sqrt{1-\zeta^2}}$ (r\'egime sous-amorti)
  \item $T_{\text{r\'eg}} \approx \pi/\omega_{co}$, \quad $T_0^{BF} \approx 2\pi/\omega_{co}$
\end{itemize}
\end{tcolorbox}

\begin{tcolorbox}[colback=ch6teal!5,colframe=ch6teal!80!black,title={R\'egulateur I pur \& double int\'egrateur en BO}]
\textbf{I pur} : $G_c(s) = K_i/s$ $\Rightarrow$ ajoute un p\^ole en 0 (type $+1$), phase BO $-90°$ suppl. $\Rightarrow$ \textbf{r\'eduit $\varphi_m$}.\\[2pt]
\textbf{Type 2} ($1/s^2$ en BO) : \'elimine $e_\infty$ pour rampe, mais tr\`es difficile \`a stabiliser. Bode : pente $-40$ dB/d\'ec, phase $= -180°$ (constant). \\[2pt]
\piege{Type\,$\geq 2$ : la BO part \`a $-180°$ d\`es les BF $\Rightarrow$ $\varphi_m$ tr\`es difficile \`a obtenir. V\'erifier par le crit\`ere du revers / les marges.}
\end{tcolorbox}

\begin{tcolorbox}[colback=ch6teal!5,colframe=ch6teal!80!black,title={Interpr\'etation plan de Nyquist}]
\begin{center}
\begin{tikzpicture}[scale=0.78, transform shape, >=Stealth,
   every node/.style={font=\fontsize{5.6}{6.4}\selectfont}]
  % --- axes ---
  \draw[->] (-2.5,0) -- (1.05,0) node[right]{Re};
  \draw[->] (0,-1.75) -- (0,1.55) node[above]{Im};
  % --- cercle unité : label déporté hors du lieu (haut-droite) ---
  \draw[gray!45, thin] (0,0) circle (1);
  \draw[gray!55, thin] (0.71,0.71) -- (1.02,1.06);
  \node[gray!70, right] at (1.0,1.1) {$|G_o|{=}1$};
  % --- lieu omega>0 (trait plein) : depart omega->0 a droite, croise l'axe reel
  %     negatif en omega_pi (entre -1 et 0 => marges OK), puis -> origine ---
  \draw[very thick, ch6teal!80!black]
     (2.0,-0.12) to[out=-105,in=20] (0.25,-1.2)
     to[out=200,in=-55] (-0.62,0)
     to[out=125,in=-150] (-0.30,0.5)
     to[out=30,in=170] (0.02,0.05);
  % --- conjugue omega<0 (tiret) ---
  \draw[very thick, ch6teal!80!black, dashed]
     (2.0,0.12) to[out=105,in=-20] (0.25,1.2)
     to[out=160,in=55] (-0.62,0);
  \node[ch6teal!75!black] at (1.55,-0.42) {$\omega{\to}0$};
  % --- point critique -1 (label en-dessous, hors zone marges) ---
  \filldraw[red!80!black] (-1,0) circle (1.7pt);
  \node[red!80!black, below left=-2pt and -2pt] at (-1,0) {$-1$};
  % --- omega_pi : croisement axe reel negatif (repere vers le bas) ---
  \filldraw[blue!75] (-0.62,0) circle (1.6pt);
  \draw[blue!75, thin] (-0.62,0) -- (-0.62,-1.05);
  \node[blue!75, below] at (-0.62,-1.08) {$\omega_\pi$};
  % --- marge de gain : double-fleche fine -1 <-> croisement, label AU-DESSUS ---
  \draw[<->, green!45!black, line width=0.45pt, >={Stealth[length=2.2pt]}]
       (-1,0.13) -- (-0.62,0.13);
  \draw[green!45!black, line width=0.3pt] (-0.81,0.15) -- (-0.81,0.36);
  \node[green!45!black, above] at (-0.81,0.34) {$1/A_m$};
  % --- omega_co : croisement du cercle unite (3e quadrant) ---
  \coordinate (co) at ({cos(213)},{sin(213)});
  \filldraw[orange!85!black] (co) circle (1.6pt);
  \draw[orange!85!black, thin] (co) -- ({1.62*cos(213)},{1.5*sin(213)});
  \node[orange!85!black, left] at ({1.68*cos(213)},{1.5*sin(213)}) {$\omega_{co}$};
  % --- marge de phase : rayon vers omega_co + arc 180deg -> angle de G_o(jw_co) ---
  \draw[gray!55, thin] (0,0) -- (co);
  \draw[red!70!black, line width=0.9pt] (180:0.58) arc (180:213:0.58);
  \draw[red!70!black, line width=0.3pt] ({0.58*cos(196.5)},{0.58*sin(196.5)}) -- (0.36,-0.54);
  \node[red!70!black, right] at (0.34,-0.57) {$\varphi_m$};
\end{tikzpicture}
\end{center}
\vars{Lieu de $G_o(\jw)$. Trait plein : $\omega>0$ ; tiret : $\omega<0$ (sym. conj.). $A_m{=}1/|G_o(\jw_\pi)|$ (croisement axe r\'eel n\'eg.). $\varphi_m$ = angle entre $-$axe r\'eel et $G_o(\jw_{co})$ sur le cercle unit\'e. Stable (revers) : $(-1,0)$ \`a gauche du lieu.}
\par\smallskip
\textbf{Propri\'et\'e marge de phase r\'egulateur} :
\[
\varphi_m(\mathrm{PI}) < \varphi_m(\mathrm{P}) < \varphi_m(\mathrm{PID}) < \varphi_m(\mathrm{PD})
\]
\end{tcolorbox}

\begin{tcolorbox}[colback=ch6teal!5,colframe=ch6teal!80!black,title={Synth\`ese fr\'equentielle : compensation p\^ole-z\'ero (ch.8)}]
\textbf{Principe} : placer un z\'ero de $G_c$ sur un p\^ole (dominant) de $G_a$ $\Rightarrow$ il dispara\^it de $G_o{=}G_cG_a$. PI : $T_i{=}\tau_{a,\max}$ ; PID : 2 z\'eros $\to$ 2 p\^oles dom. ($T_i{=}\tau_1{+}\tau_2$, $T_d{=}\tau_1\tau_2/T_i$).
\piege{Valable si $G_a$ \textbf{stable en BO}. Le p\^ole compens\'e reste dans $G_{yv}$ (maintien).}
\par\smallskip
\textbf{Deux gains \`a ne pas confondre} ($K_a$ = gain du proc\'ed\'e $G_a$) :
\tightmath{
\boxed{K_o = \frac{K_p\,K_a}{T_i}} \;(\text{gain de \textbf{boucle}}) \qquad \boxed{K_p = \frac{K_o\,T_i}{K_a}} \;(\text{gain \textbf{r\'egulateur}})
}
\textbf{M\'ethode} : 1. compenser les p\^oles dom. $\Rightarrow G_o = \frac{K_o}{s}\cdot\frac{1}{1{+}\tau_r s}$. 2. n\'egliger $\tau_r$ $\Rightarrow$ BF 1\textsuperscript{er} ordre, p\^ole en $-K_o$ : $T_{\text{r\'eg},5\%} = 3/K_o$. 3. $K_o = 3/T_{\text{r\'eg}}$ (ou lu sur Bode). 4. $K_p = K_o T_i/K_a$.
\vars{V\'erifier $\omega_{co}{=}K_o \ll 1/\tau_r$, sinon l'approx. 1\textsuperscript{er} ordre tombe.}
\par\smallskip
\textbf{Ex. PI} : $G_a = \frac{2}{(1{+}0.5s)(1{+}0.05s)}$, vise $T_{\text{r\'eg}} = 0.3$ s.
\tightmath{
T_i = \tau_{a,\max} = 0.5,\ \ K_o = \tfrac{3}{0.3} = 10 \;\Rightarrow\; K_p = \frac{K_o T_i}{K_a} = \frac{10\cdot 0.5}{2} = 2.5
}
\textbf{Ex. PID} : $G_a = \frac{K_a}{(1{+}0.5s)(1{+}0.1s)(1{+}0.02s)}$ : $T_i = 0.6$ s, $T_d = \frac{\tau_1\tau_2}{T_i} = 0.083$ s, puis $G_o = \frac{K_o}{s(1{+}0.02s)}$, $K_o = 3/T_{\text{r\'eg}}$, $K_p = K_o T_i/K_a$.
\end{tcolorbox}
